% frontmatter.tex -- Module 8 (From the Universe to the Atom): cover, How to use this guide,
% Contents, Module overview.
% \inquiry boxes of ch1.tex ... ch5.tex.
% "Formulas at a glance" = the formula cards of ch1.tex ... ch5.tex, same notation and the
% same badge (on the sheet / not on the sheet / combine). The de Broglie card appears in
% ch3 and ch5; it is listed once here (ch3) and twice in formulasheet.tex.
% Experiment list: the Key experiment boxes of ch1.tex ... ch5.tex.
%
% Local layout fixes (common/ is shared, so they are made here, inside groups):
%  - cover: break the title after "Universe" (the automatic break left "Atom" alone on
%    the last line); \moduletitle is changed for the cover only.
%  - contents: three-digit page numbers were 2 pt wider than the page-number box
%    (titletoc \@pnumwidth = 1.55em), so they stuck out into the margin. The box is made
%    1.9em wide; the correction 0.65em keeps the text margin at 2.55em (= \@tocrmarg).
\begingroup
  \def\moduletitle{From the Universe\\to the Atom}%
  \makecover
\endgroup
\IfFileExists{../common/howtouse.tex}{% howtouse.tex -- shared "How to use this guide" pages (same in all four books).
% Input it from frontmatter.tex:   % howtouse.tex -- shared "How to use this guide" pages (same in all four books).
% Input it from frontmatter.tex:   % howtouse.tex -- shared "How to use this guide" pages (same in all four books).
% Input it from frontmatter.tex:   \input{../common/howtouse}
% Numbers in the mini worked example: R = 12 / 0.50 = 24 ohm (checked in Python).
\unnumberedchapter[\faInfoCircle]{How to use this guide}

This guide explains Module~\modulenumber\ (\moduletitle) from the very beginning.
It uses simple English, short sentences and many examples. Every calculation is
shown step by step.

The chapters follow the inquiry questions of the NSW Physics Stage 6 Syllabus
(2017). At the back of the book you will find a formula summary, a glossary
(with Chinese translations), a practice set, full answers and a last-minute
summary. In the PDF, you can click any entry in the contents, any
cross-reference and any question number.

\section*{What the coloured boxes mean}

Each kind of box has its own colour and icon. A colour always means the same
thing, in every book of this series.

\begin{sidebyside}[0.5]
  \begin{bigidea}[before skip=0pt, after skip=0pt, equal height group=legend1, fontupper=\normalsize]
    The main point of a topic in two or three sentences. Read these first.
  \end{bigidea}
\nextside
  \begin{keywords}[before skip=0pt, after skip=0pt, equal height group=legend1]
    \kw{Term}{what the word means, in plain words.}
    New words are also printed in \term{bold} the first time they appear.
  \end{keywords}
\end{sidebyside}

\begin{sidebyside}[0.5]
  \begin{mistakes}[before skip=0pt, after skip=0pt, equal height group=legend2]
    Errors that students often make, and how to avoid them:
    \xmark\ wrong \enspace \cmark\ right.
  \end{mistakes}
\nextside
  \begin{examtip}[before skip=0pt, after skip=0pt, equal height group=legend2]
    How the HSC asks about this topic, and what you must write to get full
    marks.
  \end{examtip}
\end{sidebyside}

\begin{sidebyside}[0.5]
  \begin{toolkit}[before skip=0pt, after skip=0pt, equal height group=legend3]{}
    A short maths skill that you need for the next part, for example
    rearranging a formula.
  \end{toolkit}
\nextside
  \begin{refresher}[before skip=0pt, after skip=0pt, equal height group=legend3]{}
    A short reminder of a Year~11 idea that you need, for example Newton's
    laws.
  \end{refresher}
\end{sidebyside}

\begin{sidebyside}[0.5]
  \begin{tcolorbox}[hsc minor=hscqc, unbreakable, before skip=0pt, after skip=0pt,
      equal height group=legend4, title={\hscboxtitle{\faQuestionCircle}{Quick check}{}}]
    Short questions to test yourself, numbered like \textcolor{hscqc}{\bfseries Q2.3}
    (Chapter~2, question~3). The answers at the back use the same numbers.
  \end{tcolorbox}
\nextside
  \begin{experiment}[unbreakable, before skip=0pt, after skip=0pt, equal height group=legend4]{}
    A practical or historical experiment: aim, apparatus, method, results,
    conclusion and why it matters.
  \end{experiment}
\end{sidebyside}

\begin{formulacard}[unbreakable, before skip=6pt]{Name of the formula}
  \cardformula{V = IR}
  Every formula has a card: the formula, the meaning and SI unit of each
  symbol, when to use it, and a common trap. The badge under the formula tells
  you about the official HSC formula sheet:
  \begin{itemize}
    \item \hscbadgeonsheet\ \ it is printed on the sheet, so you are given it
          in the exam;
    \item \hscbadgenotonsheet\ \ you must learn it, or know how to get it;
    \item \hscbadgecombine{two formulas}\ \ it is two formulas from the sheet
          joined together.
  \end{itemize}
\end{formulacard}

\begin{workedexample}[nonumber, unbreakable]{How every worked example is set out}
  A \qty{12}{V} battery drives a current of \qty{0.50}{A} through a resistor.
  Find the resistance. (Cover everything below the \emph{Solution} line with a
  sheet of paper and try it first.)
  \given{$V = \qty{12}{V}$, \ $I = \qty{0.50}{A}$}
  \find{the resistance $R$}
  \formula{$V = IR$, \ so $R = \dfrac{V}{I}$}
  \substitute{$R = \dfrac{\qty{12}{V}}{\qty{0.50}{A}}$ \ (numbers with their units)}
  \answer{\final{$R = \qty{24}{\ohm}$} \ (with a unit and sensible significant
    figures)}
  \checkstep{is the answer sensible? Units: $\unit{V} \div \unit{A} = \unit{\ohm}$ \cmark}
\end{workedexample}

Each chapter starts with its \textbf{inquiry question} from the syllabus, a
list of \textbf{goals} and the \textbf{syllabus dot points} it covers. It ends
with a short \textbf{summary}.

\section*{How to study with this guide in the last two weeks}

The HSC Physics exam is on \textbf{Thursday 5 November 2026}. You have four
books (Modules 5 to 8). Here is a simple plan.

\begin{table}
  \centering
  \begin{tabular}{@{}L{30mm}L{\dimexpr\linewidth-34mm\relax}@{}}
    \toprule
    \colhead{When} & \colhead{What to do}\\
    \midrule
    Days 1 to 10 & Work through the four books, about two or three days for
      each book. For every chapter, follow steps 1 to 5 below.\\
    Days 11 to 13 & Do the Practice set of each book like a real exam: time
      yourself and use only the formula sheet. Mark your answers with the
      marking criteria. Redo every question you got wrong.\\
    The night before & Read the Last-minute summary of each book and look at
      the formula sheet once more. Then go to bed early.\\
    \bottomrule
  \end{tabular}
\end{table}

For each chapter:
\begin{enumerate}
  \item Read the \textbf{big idea} and the \textbf{key words} (about 10 minutes).
  \item Read the explanation and the \textbf{formula cards}. Say each formula
        ``in words''.
  \item \textbf{Cover the solution} of each worked example and try it yourself.
        Then compare your work with the book, step by step.
  \item Do the \textbf{Quick check} questions. Mark them with the answers at the
        back.
  \item Write every mistake you made in a ``mistakes list''. Read this list
        again before the exam.
\end{enumerate}

\section*{How to use the official formula sheet in the exam}

In the exam you get the NESA data sheet and formulae sheet. It gives many
formulas and constants, but not all of them.

\begin{enumerate}
  \item \textbf{Before the exam}, learn where each formula is on the sheet. The
        sheet has four parts: motion, forces and gravity; waves and
        thermodynamics; electricity and magnetism; quantum, special relativity
        and nuclear. Then you can find a formula quickly.
  \item \textbf{Learn the formulas that are not on the sheet.} Their formula
        cards have the red badge and show you how to derive them.
  \item \textbf{Check what each symbol means} in the question. The same letter
        can mean different things: for example, $T$ is a period in one formula
        and a temperature in another.
  \item \textbf{Change every quantity to SI units first}: km to m, nm to m,
        \si{km.h^{-1}} to \si{m.s^{-1}}, eV to J.
  \item \textbf{Use the values on the data sheet}, for example
        $g = \qty{9.8}{m.s^{-2}}$ and $c = \qty{3.00e8}{m.s^{-1}}$.
  \item \textbf{Show your working} in the same order as the worked examples:
        formula, substitution, answer with a unit.
\end{enumerate}

% Numbers in the mini worked example: R = 12 / 0.50 = 24 ohm (checked in Python).
\unnumberedchapter[\faInfoCircle]{How to use this guide}

This guide explains Module~\modulenumber\ (\moduletitle) from the very beginning.
It uses simple English, short sentences and many examples. Every calculation is
shown step by step.

The chapters follow the inquiry questions of the NSW Physics Stage 6 Syllabus
(2017). At the back of the book you will find a formula summary, a glossary
(with Chinese translations), a practice set, full answers and a last-minute
summary. In the PDF, you can click any entry in the contents, any
cross-reference and any question number.

\section*{What the coloured boxes mean}

Each kind of box has its own colour and icon. A colour always means the same
thing, in every book of this series.

\begin{sidebyside}[0.5]
  \begin{bigidea}[before skip=0pt, after skip=0pt, equal height group=legend1, fontupper=\normalsize]
    The main point of a topic in two or three sentences. Read these first.
  \end{bigidea}
\nextside
  \begin{keywords}[before skip=0pt, after skip=0pt, equal height group=legend1]
    \kw{Term}{what the word means, in plain words.}
    New words are also printed in \term{bold} the first time they appear.
  \end{keywords}
\end{sidebyside}

\begin{sidebyside}[0.5]
  \begin{mistakes}[before skip=0pt, after skip=0pt, equal height group=legend2]
    Errors that students often make, and how to avoid them:
    \xmark\ wrong \enspace \cmark\ right.
  \end{mistakes}
\nextside
  \begin{examtip}[before skip=0pt, after skip=0pt, equal height group=legend2]
    How the HSC asks about this topic, and what you must write to get full
    marks.
  \end{examtip}
\end{sidebyside}

\begin{sidebyside}[0.5]
  \begin{toolkit}[before skip=0pt, after skip=0pt, equal height group=legend3]{}
    A short maths skill that you need for the next part, for example
    rearranging a formula.
  \end{toolkit}
\nextside
  \begin{refresher}[before skip=0pt, after skip=0pt, equal height group=legend3]{}
    A short reminder of a Year~11 idea that you need, for example Newton's
    laws.
  \end{refresher}
\end{sidebyside}

\begin{sidebyside}[0.5]
  \begin{tcolorbox}[hsc minor=hscqc, unbreakable, before skip=0pt, after skip=0pt,
      equal height group=legend4, title={\hscboxtitle{\faQuestionCircle}{Quick check}{}}]
    Short questions to test yourself, numbered like \textcolor{hscqc}{\bfseries Q2.3}
    (Chapter~2, question~3). The answers at the back use the same numbers.
  \end{tcolorbox}
\nextside
  \begin{experiment}[unbreakable, before skip=0pt, after skip=0pt, equal height group=legend4]{}
    A practical or historical experiment: aim, apparatus, method, results,
    conclusion and why it matters.
  \end{experiment}
\end{sidebyside}

\begin{formulacard}[unbreakable, before skip=6pt]{Name of the formula}
  \cardformula{V = IR}
  Every formula has a card: the formula, the meaning and SI unit of each
  symbol, when to use it, and a common trap. The badge under the formula tells
  you about the official HSC formula sheet:
  \begin{itemize}
    \item \hscbadgeonsheet\ \ it is printed on the sheet, so you are given it
          in the exam;
    \item \hscbadgenotonsheet\ \ you must learn it, or know how to get it;
    \item \hscbadgecombine{two formulas}\ \ it is two formulas from the sheet
          joined together.
  \end{itemize}
\end{formulacard}

\begin{workedexample}[nonumber, unbreakable]{How every worked example is set out}
  A \qty{12}{V} battery drives a current of \qty{0.50}{A} through a resistor.
  Find the resistance. (Cover everything below the \emph{Solution} line with a
  sheet of paper and try it first.)
  \given{$V = \qty{12}{V}$, \ $I = \qty{0.50}{A}$}
  \find{the resistance $R$}
  \formula{$V = IR$, \ so $R = \dfrac{V}{I}$}
  \substitute{$R = \dfrac{\qty{12}{V}}{\qty{0.50}{A}}$ \ (numbers with their units)}
  \answer{\final{$R = \qty{24}{\ohm}$} \ (with a unit and sensible significant
    figures)}
  \checkstep{is the answer sensible? Units: $\unit{V} \div \unit{A} = \unit{\ohm}$ \cmark}
\end{workedexample}

Each chapter starts with its \textbf{inquiry question} from the syllabus, a
list of \textbf{goals} and the \textbf{syllabus dot points} it covers. It ends
with a short \textbf{summary}.

\section*{How to study with this guide in the last two weeks}

The HSC Physics exam is on \textbf{Thursday 5 November 2026}. You have four
books (Modules 5 to 8). Here is a simple plan.

\begin{table}
  \centering
  \begin{tabular}{@{}L{30mm}L{\dimexpr\linewidth-34mm\relax}@{}}
    \toprule
    \colhead{When} & \colhead{What to do}\\
    \midrule
    Days 1 to 10 & Work through the four books, about two or three days for
      each book. For every chapter, follow steps 1 to 5 below.\\
    Days 11 to 13 & Do the Practice set of each book like a real exam: time
      yourself and use only the formula sheet. Mark your answers with the
      marking criteria. Redo every question you got wrong.\\
    The night before & Read the Last-minute summary of each book and look at
      the formula sheet once more. Then go to bed early.\\
    \bottomrule
  \end{tabular}
\end{table}

For each chapter:
\begin{enumerate}
  \item Read the \textbf{big idea} and the \textbf{key words} (about 10 minutes).
  \item Read the explanation and the \textbf{formula cards}. Say each formula
        ``in words''.
  \item \textbf{Cover the solution} of each worked example and try it yourself.
        Then compare your work with the book, step by step.
  \item Do the \textbf{Quick check} questions. Mark them with the answers at the
        back.
  \item Write every mistake you made in a ``mistakes list''. Read this list
        again before the exam.
\end{enumerate}

\section*{How to use the official formula sheet in the exam}

In the exam you get the NESA data sheet and formulae sheet. It gives many
formulas and constants, but not all of them.

\begin{enumerate}
  \item \textbf{Before the exam}, learn where each formula is on the sheet. The
        sheet has four parts: motion, forces and gravity; waves and
        thermodynamics; electricity and magnetism; quantum, special relativity
        and nuclear. Then you can find a formula quickly.
  \item \textbf{Learn the formulas that are not on the sheet.} Their formula
        cards have the red badge and show you how to derive them.
  \item \textbf{Check what each symbol means} in the question. The same letter
        can mean different things: for example, $T$ is a period in one formula
        and a temperature in another.
  \item \textbf{Change every quantity to SI units first}: km to m, nm to m,
        \si{km.h^{-1}} to \si{m.s^{-1}}, eV to J.
  \item \textbf{Use the values on the data sheet}, for example
        $g = \qty{9.8}{m.s^{-2}}$ and $c = \qty{3.00e8}{m.s^{-1}}$.
  \item \textbf{Show your working} in the same order as the worked examples:
        formula, substitution, answer with a unit.
\end{enumerate}

% Numbers in the mini worked example: R = 12 / 0.50 = 24 ohm (checked in Python).
\unnumberedchapter[\faInfoCircle]{How to use this guide}

This guide explains Module~\modulenumber\ (\moduletitle) from the very beginning.
It uses simple English, short sentences and many examples. Every calculation is
shown step by step.

The chapters follow the inquiry questions of the NSW Physics Stage 6 Syllabus
(2017). At the back of the book you will find a formula summary, a glossary
(with Chinese translations), a practice set, full answers and a last-minute
summary. In the PDF, you can click any entry in the contents, any
cross-reference and any question number.

\section*{What the coloured boxes mean}

Each kind of box has its own colour and icon. A colour always means the same
thing, in every book of this series.

\begin{sidebyside}[0.5]
  \begin{bigidea}[before skip=0pt, after skip=0pt, equal height group=legend1, fontupper=\normalsize]
    The main point of a topic in two or three sentences. Read these first.
  \end{bigidea}
\nextside
  \begin{keywords}[before skip=0pt, after skip=0pt, equal height group=legend1]
    \kw{Term}{what the word means, in plain words.}
    New words are also printed in \term{bold} the first time they appear.
  \end{keywords}
\end{sidebyside}

\begin{sidebyside}[0.5]
  \begin{mistakes}[before skip=0pt, after skip=0pt, equal height group=legend2]
    Errors that students often make, and how to avoid them:
    \xmark\ wrong \enspace \cmark\ right.
  \end{mistakes}
\nextside
  \begin{examtip}[before skip=0pt, after skip=0pt, equal height group=legend2]
    How the HSC asks about this topic, and what you must write to get full
    marks.
  \end{examtip}
\end{sidebyside}

\begin{sidebyside}[0.5]
  \begin{toolkit}[before skip=0pt, after skip=0pt, equal height group=legend3]{}
    A short maths skill that you need for the next part, for example
    rearranging a formula.
  \end{toolkit}
\nextside
  \begin{refresher}[before skip=0pt, after skip=0pt, equal height group=legend3]{}
    A short reminder of a Year~11 idea that you need, for example Newton's
    laws.
  \end{refresher}
\end{sidebyside}

\begin{sidebyside}[0.5]
  \begin{tcolorbox}[hsc minor=hscqc, unbreakable, before skip=0pt, after skip=0pt,
      equal height group=legend4, title={\hscboxtitle{\faQuestionCircle}{Quick check}{}}]
    Short questions to test yourself, numbered like \textcolor{hscqc}{\bfseries Q2.3}
    (Chapter~2, question~3). The answers at the back use the same numbers.
  \end{tcolorbox}
\nextside
  \begin{experiment}[unbreakable, before skip=0pt, after skip=0pt, equal height group=legend4]{}
    A practical or historical experiment: aim, apparatus, method, results,
    conclusion and why it matters.
  \end{experiment}
\end{sidebyside}

\begin{formulacard}[unbreakable, before skip=6pt]{Name of the formula}
  \cardformula{V = IR}
  Every formula has a card: the formula, the meaning and SI unit of each
  symbol, when to use it, and a common trap. The badge under the formula tells
  you about the official HSC formula sheet:
  \begin{itemize}
    \item \hscbadgeonsheet\ \ it is printed on the sheet, so you are given it
          in the exam;
    \item \hscbadgenotonsheet\ \ you must learn it, or know how to get it;
    \item \hscbadgecombine{two formulas}\ \ it is two formulas from the sheet
          joined together.
  \end{itemize}
\end{formulacard}

\begin{workedexample}[nonumber, unbreakable]{How every worked example is set out}
  A \qty{12}{V} battery drives a current of \qty{0.50}{A} through a resistor.
  Find the resistance. (Cover everything below the \emph{Solution} line with a
  sheet of paper and try it first.)
  \given{$V = \qty{12}{V}$, \ $I = \qty{0.50}{A}$}
  \find{the resistance $R$}
  \formula{$V = IR$, \ so $R = \dfrac{V}{I}$}
  \substitute{$R = \dfrac{\qty{12}{V}}{\qty{0.50}{A}}$ \ (numbers with their units)}
  \answer{\final{$R = \qty{24}{\ohm}$} \ (with a unit and sensible significant
    figures)}
  \checkstep{is the answer sensible? Units: $\unit{V} \div \unit{A} = \unit{\ohm}$ \cmark}
\end{workedexample}

Each chapter starts with its \textbf{inquiry question} from the syllabus, a
list of \textbf{goals} and the \textbf{syllabus dot points} it covers. It ends
with a short \textbf{summary}.

\section*{How to study with this guide in the last two weeks}

The HSC Physics exam is on \textbf{Thursday 5 November 2026}. You have four
books (Modules 5 to 8). Here is a simple plan.

\begin{table}
  \centering
  \begin{tabular}{@{}L{30mm}L{\dimexpr\linewidth-34mm\relax}@{}}
    \toprule
    \colhead{When} & \colhead{What to do}\\
    \midrule
    Days 1 to 10 & Work through the four books, about two or three days for
      each book. For every chapter, follow steps 1 to 5 below.\\
    Days 11 to 13 & Do the Practice set of each book like a real exam: time
      yourself and use only the formula sheet. Mark your answers with the
      marking criteria. Redo every question you got wrong.\\
    The night before & Read the Last-minute summary of each book and look at
      the formula sheet once more. Then go to bed early.\\
    \bottomrule
  \end{tabular}
\end{table}

For each chapter:
\begin{enumerate}
  \item Read the \textbf{big idea} and the \textbf{key words} (about 10 minutes).
  \item Read the explanation and the \textbf{formula cards}. Say each formula
        ``in words''.
  \item \textbf{Cover the solution} of each worked example and try it yourself.
        Then compare your work with the book, step by step.
  \item Do the \textbf{Quick check} questions. Mark them with the answers at the
        back.
  \item Write every mistake you made in a ``mistakes list''. Read this list
        again before the exam.
\end{enumerate}

\section*{How to use the official formula sheet in the exam}

In the exam you get the NESA data sheet and formulae sheet. It gives many
formulas and constants, but not all of them.

\begin{enumerate}
  \item \textbf{Before the exam}, learn where each formula is on the sheet. The
        sheet has four parts: motion, forces and gravity; waves and
        thermodynamics; electricity and magnetism; quantum, special relativity
        and nuclear. Then you can find a formula quickly.
  \item \textbf{Learn the formulas that are not on the sheet.} Their formula
        cards have the red badge and show you how to derive them.
  \item \textbf{Check what each symbol means} in the question. The same letter
        can mean different things: for example, $T$ is a period in one formula
        and a temperature in another.
  \item \textbf{Change every quantity to SI units first}: km to m, nm to m,
        \si{km.h^{-1}} to \si{m.s^{-1}}, eV to J.
  \item \textbf{Use the values on the data sheet}, for example
        $g = \qty{9.8}{m.s^{-2}}$ and $c = \qty{3.00e8}{m.s^{-1}}$.
  \item \textbf{Show your working} in the same order as the worked examples:
        formula, substitution, answer with a unit.
\end{enumerate}
}{}
\begingroup
  \contentsmargin[0.65em]{1.9em}%
  \tableofcontents
\endgroup

\unnumberedchapter[\faMap]{Module overview}

% ---- helpers for the overview only (prefix ov)
% sheet marks: on the sheet / learn it / combine (same colours and icons as \frow)
\newcommand\ovYes{\textcolor{hscformula}{\faCheck}}
\newcommand\ovNo{\textcolor{hscmistake}{\faTimes}}
\newcommand\ovComb{\textcolor{hsctiptext}{\faLink}}
% one row of the formula grid:  \ovrow{mark}{formula}{what it is for}
\newcommand\ovrow[3]{#1 & #2 & {\footnotesize #3}\\\hline}
% chapter heading row of the formula grid
\newcommand\ovch[2]{\multicolumn{3}{@{}l@{}}{\rule{0pt}{3.4ex}\cref{#1}:
  {\bfseries\color{hscbrand}#2}}\\\hline}
\newenvironment{ovgrid}{%
  \small\renewcommand\arraystretch{1.4}\arrayrulecolor{hscrule}%
  \begin{tabular}[t]{@{}C{4.5mm}@{\hspace{1.5mm}}L{40mm}@{\hspace{1.5mm}}L{\dimexpr\linewidth-47.5mm\relax}@{}}}%
  {\end{tabular}}

\section*{What this module is about}

Module~8 asks two big questions. Where did the elements come from? And what is an atom made
of? It starts with the whole Universe: the Big Bang, the light from galaxies, and the stars
that build new elements. Then it zooms in: the electron, the nucleus, quantum ideas,
radioactivity, fission and fusion, and last the quarks inside the proton.

Much of this module is about \emph{evidence}. For each experiment, know what was done, what
was seen, and what it showed.

\section*{The inquiry questions}

The syllabus asks five inquiry questions (big questions). Each chapter answers one of them
(\cref{ov:tab:iq}).

\begin{table}
  \caption{The five inquiry questions of Module~8 and where this book answers them.}
  \label{ov:tab:iq}
  \centering
  \small
  \begin{tabular}{@{}L{62mm}L{17mm}L{60mm}@{}}
    \toprule
    \colhead{Inquiry question} & \colhead{Chapter} & \colhead{What the chapter answers}\\
    \midrule
    \textbf{IQ1}: \emph{What evidence is there for the origins of the elements?}
      & \cref{ch1:chap}
      & The Big Bang, Hubble's law, spectra, the HR diagram, fusion in stars.\\
    \midrule
    \textbf{IQ2}: \emph{How is it known that atoms are made up of protons, neutrons and
      electrons?}
      & \cref{ch2:chap}
      & Thomson's electron, Millikan's charge, Rutherford's nucleus, Chadwick's neutron.\\
    \midrule
    \textbf{IQ3}: \emph{How is it known that classical physics cannot explain the
      properties of the atom?}
      & \cref{ch3:chap}
      & Bohr's energy levels, hydrogen's spectrum, matter waves, Schr\"odinger.\\
    \midrule
    \textbf{IQ4}: \emph{How can the energy of the atomic nucleus be harnessed?}
      & \cref{ch4:chap}
      & Radioactive decay, half-life, binding energy, fission and fusion.\\
    \midrule
    \textbf{IQ5}: \emph{How is it known that human understanding of matter is still
      incomplete?}
      & \cref{ch5:chap}
      & Quarks, the Standard Model, particle accelerators, open questions.\\
    \bottomrule
  \end{tabular}
\end{table}

\section*{How the parts fit together}

In \cref{ov:fig:map}, each arrow means ``this idea is used in the next one''.

\begin{figure}
  \centering
  % layout: top row Ch 2 -> Ch 3 -> Ch 1; bottom row Ch 5 <- Ch 4, each arrow on its own
  % (no crossings); boxes 33 mm wide, gaps 2.2 cm for the labels; total width 14.3 cm.
  \begin{tikzpicture}[x=1cm, y=1cm,
      ch/.style={concept=29mm, minimum height=19mm}]
    \node[concept=70mm, fill=hsctool!8, draw=hsctool!60] (tools) at (-2.75,0)
      {\textbf{From Modules 6 and 7:} forces on charges in fields; photons
       (\mbox{$E = hf$}); \mbox{$E = mc^2$}};
    \node[ch] (c2) at (-5.5,-2.35) {\textbf{Ch 2}\\ Structure of the atom: electron,
      nucleus, neutron};
    \node[ch] (c3) at (0,-2.35) {\textbf{Ch 3}\\ Quantum atom: energy levels, matter
      waves};
    \node[ch] (c1) at (5.5,-2.35) {\textbf{Ch 1}\\ Origins of the elements: Big Bang and
      stars};
    \node[ch] (c5) at (-5.5,-5.15) {\textbf{Ch 5}\\ Deep inside the atom: quarks,
      accelerators};
    \node[ch] (c4) at (0,-5.15) {\textbf{Ch 4}\\ The nucleus: decay, fission and fusion};
    \draw[link] (tools.south -| c2.north) -- node[linklabel, right=2pt] {fields bend beams}
      (c2.north);
    \draw[link] (tools.south -| c3.north) -- node[linklabel, right=2pt] {photons} (c3.north);
    \draw[link] (c2.east) -- node[linklabel, above=2pt, align=center]
      {Rutherford's\\atom fails} (c3.west);
    \draw[link] (c3.east) -- node[linklabel, above=2pt, align=center]
      {levels explain\\star spectra} (c1.west);
    \draw[link] ($(c2.south)+(-0.4,0)$) -- node[linklabel, right=2pt]
      {inside $\mathrm{p}$ and $\mathrm{n}$?} ($(c5.north)+(-0.4,0)$);
    \draw[link] ($(c2.south east)+(-0.1,0)$) -- node[linklabel, pos=0.55, right=3pt]
      {the nucleus} ($(c4.north west)+(0.2,0)$);
    \draw[link] (c4.west) -- node[linklabel, below=3pt] {beta decay} (c5.east);
    \draw[link] ($(c4.north east)+(-0.2,0)$) -- node[linklabel, pos=0.5, below right=1pt]
      {fusion powers stars} ($(c1.south west)+(0.3,0)$);
  \end{tikzpicture}
  \caption{How the ideas in this module are linked (``Ch'' means chapter).}\label{ov:fig:map}
\end{figure}

\section*{Key experiments}

Each of these experiments has a Key experiment box in its chapter
(\cref{ov:tab:experiments}). Learn what was done, what was seen and what it showed.

\begin{table}
  \caption{Key experiments. A date means a historical experiment.}\label{ov:tab:experiments}
  \centering
  \small
  \begin{tabular}{@{}L{64mm}L{56mm}L{18mm}@{}}
    \toprule
    \colhead{Experiment} & \colhead{What it shows} & \colhead{Where}\\
    \midrule
    Hubble's law from galaxy data & redshift grows with distance & \cref{ch1:sec:experiments}\\
    Three kinds of spectra (practical) & hot solids: all colours; gases: lines
      & \cref{ch1:sec:experiments}\\
    Plotting an HR diagram from star data & main sequence, giants, white dwarfs
      & \cref{ch1:sec:experiments}\\
    \midrule
    Cathode-ray experiments (1870s--1890s) & negative particles, bent by fields
      & \cref{ch2:sec:experiments}\\
    Thomson's charge-to-mass ratio (1897) & electrons: far lighter than atoms
      & \cref{ch2:sec:experiments}\\
    Millikan's oil drops (1909--1913) & charge is quantised, in steps of $e$
      & \cref{ch2:sec:experiments}\\
    Geiger--Marsden gold foil (1909) & a tiny, dense, positive nucleus
      & \cref{ch2:sec:experiments}\\
    Chadwick's neutron (1932) & neutral, with about a proton's mass
      & \cref{ch2:sec:experiments}\\
    \midrule
    Balmer series of hydrogen (practical) & the lines fit the Rydberg equation
      & \cref{ch3:sec:experiments}\\
    Davisson and Germer (1927) & electrons diffract, with $\lambda = \frac{h}{mv}$
      & \cref{ch3:sec:experiments}\\
    G.\,P. Thomson's electron rings (1927) & electrons diffract through thin foil
      & \cref{ch3:sec:experiments}\\
    \midrule
    Half-life with dice (practical) & random decay, steady half-life
      & \cref{ch4:sec:experiments}\\
    Penetrating power of $\alpha$, $\beta$, $\gamma$ (practical)
      & $\alpha$ stopped by paper, $\beta$ by aluminium & \cref{ch4:sec:experiments}\\
    \midrule
    Deep inelastic scattering, SLAC (1967--73) & protons contain point-like quarks
      & \cref{ch5:sec:experiments}\\
    Higgs boson at the LHC (2012) & completes the Standard Model
      & \cref{ch5:sec:experiments}\\
    \bottomrule
  \end{tabular}
\end{table}

The text also explains Rutherford's first transmutation (1919) and the discovery of fission
(1938) in \cref{ch4:chap}, and Anderson's positron (1932) in \cref{ch5:chap}.

\section*{All the formulas at a glance}

Every formula card in this book. \ovYes\ = on the official HSC formula sheet.
\ovNo\ = not on the sheet: learn it or derive it. \ovComb\ = not printed in this form:
join two items from the sheet. \hyperref[fs:chap]{\emph{Formula sheet explained}}
(page~\pageref{fs:chap}) says when to use each one.

\par\smallskip\noindent
\begin{minipage}[t]{0.488\linewidth}
\begin{ovgrid}
  \ovch{ch1:chap}{Origins of the elements}
  \ovrow{\ovYes}{$E = mc^2$}{mass and energy; pairs}
  \ovrow{\ovNo}{$z = \frac{\Delta\lambda}{\lambda_0} \approx \frac{v}{c}$}{redshift of a galaxy}
  \ovrow{\ovNo}{$v = H_0 d$, \ $t \approx \frac{1}{H_0}$}{Hubble's law; age of the Universe}
  \ovrow{\ovYes}{$\lambda_\mathrm{max} = \frac{b}{T}$}{temperature of a star}
  \ovrow{\ovNo}{$\frac{L_A}{L_B} = 100^{\frac{M_B - M_A}{5}}$}{luminosities from magnitudes}
  \ovrow{\ovComb}{$E = \Delta m\,c^2$\newline $E\,(\mathrm{MeV}) = \Delta m\,(\mathrm{u})
    \times 931.5$}{energy from fusion in stars}
  \ovch{ch2:chap}{Structure of the atom}
  \ovrow{\ovNo}{$v = \frac{E}{B}$}{velocity selector}
  \ovrow{\ovNo}{$\frac{q}{m} = \frac{E}{B^2 r}$}{Thomson's $q/m$}
  \ovrow{\ovNo}{$\frac{q}{m} = \frac{2V}{B^2 r^2}$}{school e/m tube}
  \ovrow{\ovNo}{$q = \frac{mgd}{V}$}{Millikan's oil drop}
  \ovrow{\ovNo}{$q = ne$}{charge is quantised}
  \ovrow{\ovYes}{$F = \frac{1}{4\pi\varepsilon_0}\frac{q_1 q_2}{r^2}$}{Coulomb's law: alpha
    and nucleus}
  \ovrow{\ovYes}{$\sum m\vec{v}_\mathrm{before}$\newline
    $\quad = \sum m\vec{v}_\mathrm{after}$\newline
    $\sum\frac{1}{2}mv^2_\mathrm{before}$\newline
    $\quad = \sum\frac{1}{2}mv^2_\mathrm{after}$}{Chadwick's collisions}
  \ovch{ch3:chap}{The quantum atom}
  \ovrow{\ovYes}{$E = hf$}{photon energy}
  \ovrow{\ovComb}{$E = \frac{hc}{\lambda}$}{photon energy from $\lambda$}
  \ovrow{\ovYes}{$\frac{1}{\lambda} = R\left(\frac{1}{n_\mathrm{f}^2} - \frac{1}{n_\mathrm{i}^2}\right)$}
    {hydrogen's lines}
\end{ovgrid}
\end{minipage}\hfill
\begin{minipage}[t]{0.488\linewidth}
\begin{ovgrid}
  \ovch{ch3:chap}{(continued)}
  \ovrow{\ovYes}{$\lambda = \frac{h}{mv}$}{matter waves (also Ch~5)}
  \ovrow{\ovNo}{$2\pi r = n\lambda$}{allowed Bohr orbits}
  \ovch{ch4:chap}{The nucleus}
  \ovrow{\ovNo}{$\sum A_\mathrm{before} = \sum A_\mathrm{after}$\newline
    $\sum Z_\mathrm{before} = \sum Z_\mathrm{after}$}{balance a nuclear equation}
  \ovrow{\ovYes}{$N_\mathrm{t} = N_0 e^{-\lambda t}$}{amount left after $t$}
  \ovrow{\ovYes}{$\lambda = \frac{\ln 2}{t_{\frac{1}{2}}}$}{decay constant}
  \ovrow{\ovNo}{$N_\mathrm{t} = N_0\left(\frac{1}{2}\right)^{n}$, \ $n = \frac{t}{t_{\frac{1}{2}}}$}
    {whole half-lives}
  \ovrow{\ovNo}{$A = \lambda N$, \ $A_\mathrm{t} = A_0 e^{-\lambda t}$}{activity}
  \ovrow{\ovYes}{$E = mc^2$}{energy from mass lost}
  \ovrow{\ovNo}{$\Delta m = (Zm_\mathrm{p} + Nm_\mathrm{n})$\newline
    $\phantom{\Delta m =} - m_\mathrm{nucleus}$\newline
    $E_\mathrm{B} = \Delta m\,c^2$}{mass defect, binding energy}
  \ovrow{\ovNo}{$E = \sum E_\mathrm{B,after}$\newline $\phantom{E =} - \sum E_\mathrm{B,before}$}
    {energy from binding energies}
  \ovch{ch5:chap}{Deep inside the atom}
  \ovrow{\ovYes}{$p_\mathrm{v} = \frac{m_0 v}{\sqrt{\left(1 - \frac{v^2}{c^2}\right)}}$}
    {momentum near $c$}
  \ovrow{\ovNo}{$Q = q_1 + q_2 + q_3$ (baryon)\newline $Q = q_1 + q_2$ (meson)}
    {charge of a hadron}
  \ovrow{\ovYes}{$W = qV$}{energy from a gap}
  \ovrow{\ovNo}{$r = \frac{mv}{qB}$, \ $T = \frac{2\pi m}{qB}$}{circles in a magnetic field;
    cyclotron}
\end{ovgrid}
\end{minipage}
\par\medskip
